% !PRINTED 102
\weberpage{102}
%<<CONT
the case $m = 8$, at least two different primes must be contained in $m$.

Let us set:
\[
m = q_1\, q_2 \cdots q_\mu,
\]
if $m$ is odd or divisible only by 4,
\[
m = q_2\, q_3 \cdots q_\mu,
\]
if $m$ is divisible by 8, so that $q_1 = 4$ may occur and in the second formula $q_2 = 8$, while the remaining $q$ are odd primes; then the index moduli in the three cases are:
\[
\begin{aligned}
& q_1 - 1,\quad q_2 - 1,\ \ldots,\ q_\mu - 1\\
& 2\phantom{{}_1 -{}1,}\ ,\quad q_2 - 1,\ \ldots,\ q_\mu - 1\\
& 2,\quad 2\phantom{{}_1 -{}1}\ ,\quad q_3 - 1,\ \ldots,\ q_\mu - 1.
\end{aligned}
\]
Since by this all the index moduli are divisible by 2, the conditions for $m$ are exhausted.

We have [\S.\,21, (11)]:
\weq{1}{\delta_{k,h} = 2,\ i_{k,h} = 1,\ c_{k,h} = \tfrac{1}{2} c_h,\ k = 1, 2,\ h = 1, 2, \ldots \mu.}
Now the quantities $e_{k,h}$ are to be determined in such a way that from the matrix
\weq{2}{\begin{array}{cccc}
e_{1,1}, & e_{1,2}, & \ldots & e_{1,\mu}\\
e_{2,1}, & e_{2,2}, & \ldots & e_{2,\mu}
\end{array}}
at least one odd two-rowed determinant can be formed.

In order that $\mathfrak{A}$ be a primary divisor of $\mathfrak{N}$, the condition \S.\,21, 6. is further added, where it is now to be distinguished whether $m$ is divisible by 8 or not. If $m$ is divisible by 8, the number of index moduli is greater by one than the number of prime factors of $m$, and there is an index modulus to which no prime factor of $m$ corresponds; in the other case the two numbers agree. According to \S.\,21, 6. the $e_{k,h}$ are to be chosen in such a way that in none of the pairs
\[
e_{1,1}, e_{2,1};\ e_{1,2}, e_{2,2};\ \ldots\ e_{1,\mu}, e_{2,\mu},
\]
to which a prime factor of $m$ corresponds are both numbers even. If $m$ is divisible by 8, there is among them a pair corresponding to the index modulus $c_k = 2$, on which no such restriction is imposed. If the numbers $e_{k,h}$ are (with respect to the modulus 2) so determined, then
%CONT>>
